\chapter{Introduction}
\label{ch:intro}

This chapter motivates the project, states its purpose and contributions, and
outlines the structure of the remainder of the report.

\section{Background and Motivation}
\label{sec:intro-motivation}

The proliferation of recorded lectures and digital course material has created a
fundamental imbalance in self directed learning: students have more content than
ever, but fewer tools with which to extract, consolidate and assess their
understanding of it. A single semester course now routinely produces dozens of
hours of recorded lectures, slide decks and reference chapters. Reviewing that
material passively  re-watching a video, re-reading a chapter is one of the
least effective study strategies available. Decades of work in the learning
sciences instead point to \emph{retrieval practice}: repeatedly testing oneself on
the material produces substantially more durable learning than re-exposure to it.
The obstacle is supply. Producing good practice questions is slow, expert work,
and the questions that do exist are concentrated at the end of textbook chapters,
are not tied to what a particular instructor actually emphasised in class, and run
out long before a student has mastered a weak topic.

Existing AI assisted study tools address this gap only partially. Text based
chatbots require the student to manually transcribe or paste lecture content before
they can help, which puts the burden of the hardest step getting the lecture out
of the video and into text back on the student. Cloud dependent question
generation services remove that burden, but they are inaccessible in
bandwidth limited regions, they charge per token at a scale that makes routine
practice expensive, and they require sending course material  which may be
copyrighted or institutionally restricted  to a third party server. Crucially,
none of these tools close the full loop from raw instructional media, through
intelligent content understanding, to graded and personalised assessment without
human intervention at each stage.

There is also a quality problem that is easy to overlook. A large language model
asked to write a multiple choice question will almost always produce something
that \emph{reads} like a good question. Whether the option it marks as correct
actually \emph{is} correct is a separate matter, and one that fluency based
evaluation cannot detect. A study tool that confidently teaches a student the wrong
answer is worse than no tool at all. Any serious system for automatic question
generation therefore has to be evaluated on answer correctness, and on whether it
keeps producing genuinely different questions rather than paraphrasing the same one,
not merely on whether its output is well written.

\section{Purpose and Goal of the Project}
\label{sec:intro-purpose}

The goal of this project is to build and rigorously evaluate \textbf{LectureAssist},
a fully local, API free system that ingests a lecture video or an academic document
and autonomously produces validated, difficulty controlled assessment material,
together with a retrieval grounded chat interface over the same content.

The system is built around three ideas working in concert. First, a
\emph{dualchannel extraction} layer simultaneously processes the visual channel of
a lecture  using scene type adaptive OCR that distinguishes slides, blackboards,
whiteboards and on-screen text, each of which needs different image preprocessing
and its audio channel via Faster Whisper transcription, producing a unified,
timestamped lecture timeline. Second, a \emph{multi-agent generation pipeline}
decomposes quiz creation into specialist agents for planning, retrieval,
generation, validation and targeted refinement, enabling structured quality control
at the level of the individual question. Third, an \emph{adaptive difficulty engine}
maintains per topic performance histories across sessions and steers question
allocation toward a student's weak areas. The entire system runs on locally hosted
Gemma~3 models through the Ollama runtime, requires no external API keys, and is
packaged for single GPU deployment on Google Colab with Google Drive session
persistence.

The key contributions of this work are:

\begin{itemize}
  \item A unified multi modal lecture understanding pipeline that fuses
        scene adaptive OCR with filtered speech transcription into a synchronised
        content timeline, supporting both video and document inputs.

  \item A multi-agent quiz architecture implementing a Plan $\rightarrow$ Retrieve
        $\rightarrow$ Generate $\rightarrow$ Validate $\rightarrow$ Refine cycle
        across five question types (MCQ, True/False, Fill-in-the-Blank, Short Answer
        and Flashcard), with per-question difficulty, grounding and
        instruction-compliance validation, and a validator that chooses \emph{how}
        to remediate a rejected question rather than merely rejecting it.

  \item An adaptive assessment engine that tracks per-topic student performance
        across sessions and personalises question difficulty and distribution
        without manual configuration.

  \item A multi-tier grading system combining deterministic matching, fuzzy string
        similarity and LLM-based rubric evaluation for free-text short answers, with
        targeted concept-level feedback on incorrect responses.

  \item A fully local, API-free deployment on consumer hardware using the Ollama
        inference runtime, with RAG-enabled semantic chat and persistent session
        storage.

  \item A large-scale evaluation harness that compares the agentic pipeline against
        three direct-prompting baselines --- plain, chain-of-thought and
        program-of-thought --- and which, unlike prior automatic question generation
        evaluations, separates \emph{question quality} from \emph{answer
        correctness} by having a second, larger Answer Generator agent independently
        re-solve every generated question without seeing the marked option.
\end{itemize}

\section{Organization of the Report}
\label{sec:intro-org}

The remainder of this report is organised as follows. Chapter~\ref{ch:background}
provides the background needed to follow the rest of the report --- large language
models and the prompting strategies used as baselines, retrieval-augmented
generation, agentic architectures, speech recognition and OCR --- then reviews the
research literature and the existing commercial applications closest to this work,
and identifies the gaps this project addresses. Chapter~\ref{ch:methodology}
describes the system design: the extraction layer, the agent architecture, the four
question-generation pipelines with their prompts given verbatim, the two-agent
answer cross-check, and the software components used to build it.
Chapter~\ref{ch:results} presents the evaluation: the environment, the metrics and
exactly how each is computed, the large-scale multiple-choice results, and a
discussion of what they show --- including a failure mode of the agentic pipeline
that the evaluation exposed. Chapter~\ref{ch:impacts} discusses the societal,
ethical and environmental implications and the design constraints they impose.
Chapter~\ref{ch:planning} presents the project plan and budget.
Chapter~\ref{ch:cep} maps the work onto the complex engineering problem and
activity attributes. Chapter~\ref{ch:conclusion} summarises the project, states its
limitations honestly, and sets out the work that follows from it.
